% Valeo Cairo: Software Engineer - Test Automation, POWER (REQ2026081769). Prepared 16 Nov 2026.
% Every fact comes from facts.md. Build: TEXINPUTS=../..: pdflatex (twice).
\documentclass[10pt]{article}
\usepackage{cv-style}
\hypersetup{pdftitle={Abdelrahman Khalifa - Embedded Software Engineer - CV},pdfauthor={Abdelrahman Khalifa}}

\begin{document}

\cvheader{Abdelrahman Khalifa}{Software Engineer \enspace|\enspace Test Automation, Python, Embedded}

\section*{Summary}
Software engineer with a firmware background who tests what he builds: Python and C at EKSON's motion-simulator company (firmware, desktop app, CI), testing on the real hardware, and end-to-end tests with pytest. Knows ECU flashing and diagnostics from a UDS bootloader on STM32, and CI from building multi-architecture test pipelines as a Google Summer of Code 2026 contributor with The Linux Foundation (OpenPrinting).

\section*{Experience}
\cvjob{Embedded Software Engineer}{EKSON}{Aug 2025 -- Present}
\cvsub{Software for 3-DoF and 6-DoF pneumatic motion simulators: firmware, desktop app, backend and CI}
\begin{itemize}
  \item Wrote the ESP32-S3 firmware and migrated it from ESP-IDF/FreeRTOS to Zephyr RTOS for more deterministic timing, devicetree-based drivers and easier porting to other boards.
  \item Implemented closed-loop PID control of the pneumatic valves, with feedback from an MPU6050 IMU over I2C (calibration and a Mahony filter for orientation), plus safety stops and a watchdog.
  \item Developed the PyQt5 desktop app that controls the simulator over a UART protocol; cut its CPU usage and latency by about 40\% (measured) using asyncio to wait for ESP32 and server messages concurrently, off the UI thread.
  \item Engineered remote operation and firmware updates: a Node.js/Express + MongoDB backend that the desktop app polls over REST, then flashes new firmware with esptool; GitHub Actions CI.
\end{itemize}
\cvjob{Google Summer of Code Contributor}{The Linux Foundation (OpenPrinting)}{May 2026 -- Oct 2026}
\cvsub{Open-source project: Printers page for COSMIC Settings, in Rust \enspace|\enspace \cvlink{https://github.com/Abd002/GSOC-2026}{github.com/Abd002/GSOC-2026}}
\begin{itemize}
  \item Built the \textbf{Printers page for COSMIC Settings} (printer status, supply levels, print queues, adding printers) as 5 Rust crates, with the CUPS/IPP/DNS-SD backend isolated behind a varlink IPC service (\cvlink{https://github.com/pop-os/cosmic-settings/pull/2182}{cosmic-settings \#2182}, in review).
  \item Wrote Rust-to-C FFI bindings (libcups3 DNS-SD), ported cups-rs to libcups3 with CI, and fixed crashes in libcups' Avahi code; \textbf{11 PRs merged upstream} (\cvlink{https://github.com/OpenPrinting/cups-rs/pulls?q=is:pr+author:Abd002+is:merged}{cups-rs}, \cvlink{https://github.com/OpenPrinting/libcups/pull/164}{\#164}, \cvlink{https://github.com/OpenPrinting/libcups/pull/165}{\#165}).
  \item Set up the CI: a formatting check plus a build-and-test matrix on x86\_64 and aarch64 that builds libcups3 from source and runs the tests against real CUPS and Avahi services.
\end{itemize}
\cvjob{Embedded Linux Systems Intern}{Orange Digital Center (ODC)}{Dec 2024 -- Jan 2025}
\begin{itemize}
  \item Created custom Linux images with Yocto and Buildroot and wrote a Linux kernel module; worked with device trees, root file systems and hardware flashing. Scored 100\% in the final evaluation.
\end{itemize}

\section*{Open Source Projects}
\begin{itemize}
  \item \textbf{Zephyr RTOS} (Mar 2025): implemented the \texttt{IP\_MULTICAST\_LOOP} socket option for IPv4 across 6 networking modules, reviewed and merged upstream (\cvlink{https://github.com/zephyrproject-rtos/zephyr/issues/86977}{issue \#86977}).
\end{itemize}

\section*{Projects}
\begin{itemize}
  \item \textbf{AUTOSAR-Based V2X Vehicle} (graduation project, team; STM32, ESP32): owned the firmware update path. An ESP32 gateway downloads the image from Firebase Storage and sends it over UART with UDS services (request download 0x34, transfer data 0x36, transfer exit 0x37, read version); the STM32 bootloader verifies it with the hardware CRC unit and uses two banks with rollback. \cvlink{https://github.com/Abd002/Graduation-Project}{GitHub}
  \item \textbf{Function Solver and Plotter} (Python, PySide2, SymPy, Matplotlib): input validation, error handling and end-to-end tests with pytest-qt. \cvlink{https://github.com/Abd002/py-func-solver}{GitHub}
  \item \textbf{Seat Heater Control System} (Tiva C, FreeRTOS): heater controller for two car seats with event groups, periodic tasks and a mutex-protected UART display; sensor-failure diagnostics logged to memory, and run-time measurements on the target (CPU load, task execution and lock times, using GPTM). \cvlink{https://github.com/Abd002/SeatHeaterControlSystem}{GitHub}
  \item \textbf{Database Systems, CMU 15-445} (C++17): thread-safe Count-Min Sketch for the BusTub database (one mutex per row; insert, count, merge, top-k); passes the course's autograder tests, clean under clang-tidy, no memory leaks. \cvlink{https://github.com/Abd002/db-systems-journey}{GitHub}
\end{itemize}

\section*{Technical Skills}
\cvskill{Languages}{Python (asyncio, pytest), C, Modern C++, Bash, Rust}
\cvskill{Testing \& CI}{Test design and execution, unit, integration and end-to-end tests, testing on the hardware target, log analysis and root-cause debugging, GitHub Actions CI, Git}
\cvskill{Automotive \& embedded}{UDS diagnostics and flashing, CAN, UART, I2C, SPI; AUTOSAR (MCAL); STM32, ESP32-S3; Zephyr, FreeRTOS}

\section*{Education}
\cvjob{Bachelor of Science (B.Sc.), Computer and Systems Engineering}{Minia University}{2020 -- 2025}
\begin{itemize}
  \item Grade: Very Good with Honors.
  \item Activities: Codeforces Expert; 4-time ECPC finalist (2021--2024); Treasurer, IEEE Student Branch.
\end{itemize}

\end{document}
